\chapter{Risultati Sperimentali}


	\section{Transfer Learning Domain-Specific: L'Approccio SIPAKMED}
	
	Una delle sfide principali nell'addestramento di reti neurali profonde per task medici è la necessità di dataset annotati molto ampi. Nel nostro caso, il dataset HPV disponibile conta 133,704 patch estratte da 19 pazienti, con una distribuzione fortemente sbilanciata: solo 6,181 patch (4.62\%) sono HPV-positive, mentre le restanti 127,523 rappresentano tessuto normale. Sebbene questo numero di patch sia significativo, la variabilità biologica rappresentata da 19 soggetti pone il rischio concreto di overfitting su caratteristiche paziente-specifiche anziché su pattern patologici generalizzabili.
	
	Per affrontare questa criticità, abbiamo adottato una strategia a due livelli: (1) \textbf{Leave-One-Patient-Out cross-validation} su 3 pazienti principali con dataset sufficientemente ampio per consentire split bilanciati, e (2) \textbf{inclusione di 16 pazienti aggiuntivi} nel training per aumentare la variabilità morfologica e ridurre il rischio di overfitting. La Tabella~\ref{tab:dataset_patient_strategy} riassume questa composizione.
	
	\begin{table}[h]
		\centering
		\small
		\caption{Strategia di utilizzo dei pazienti nel dataset HPV}
		\label{tab:dataset_patient_strategy}
		\begin{tabular}{lccl}
			\toprule
			\textbf{Categoria} & \textbf{N Pazienti} & \textbf{Patches} & \textbf{Ruolo} \\
			\midrule
			LOPO CV (rotazione test) & 3 & 78,702 (59\%) & Test/Train/Val alternati \\
			Training augmentation & 16 & 55,002 (41\%) & Sempre in Train/Val \\
			\midrule
			\textbf{Totale} & \textbf{19} & \textbf{133,704 (100\%)} & — \\
			\bottomrule
		\end{tabular}
	\end{table}
	
	La strategia del Transfer Learning rappresenta la soluzione standard al problema della scarsità di dati annotati: inizializzare i pesi della rete con un modello pre-addestrato su un dataset molto più ampio. Tipicamente, in Computer Vision, si utilizza ImageNet, un dataset di oltre 14 milioni di immagini naturali. Tuttavia, la letteratura recente~\cite{raghu2019transfusion} ha dimostrato che per task medici, un pre-training domain-specific (cioè su immagini mediche) può fornire performance superiori rispetto a ImageNet, anche se il dataset di pre-training è di dimensioni molto inferiori.
	
	In questa sezione documentiamo un primo approccio basato su transfer learning dal dataset SIPAKMED, un corpus pubblico di immagini citologiche cervicali che condivide il dominio biologico con il nostro task di classificazione HPV. Presentiamo i risultati iniziali di questa strategia e, attraverso un'analisi critica delle criticità emerse, descriviamo le ottimizzazioni metodologiche che hanno permesso di raggiungere performance clinicamente rilevanti.
	
	%═══════════════════════════════════════════════════════════════════════════════
	\subsection{Il Dataset SIPAKMED per Pre-training}
	
	SIPAKMED~\cite{plissiti2018sipakmed} è un dataset pubblicamente disponibile composto da 4,049 immagini di cellule cervicali isolate, acquisite tramite Pap-test e colorate con la colorazione di Papanicolaou. Le immagini sono suddivise in 5 classi morfologiche:
	
	\begin{itemize}
		\item \textbf{Superficial-Intermediate cells}: Cellule epiteliali superficiali normali.
		\item \textbf{Parabasal squamous epithelial cells}: Cellule basali dell'epitelio squamoso.
		\item \textbf{Koilocytotic cells}: Cellule infette da HPV, caratterizzate dal tipico alone perinucleare (coilocitosi).
		\item \textbf{Metaplastic cells}: Cellule metaplastiche benigne.
		\item \textbf{Dyskeratotic cells}: Cellule con cheratinizzazione anomala.
	\end{itemize}
	
	La presenza della classe \textit{Koilocytotic cells} rende questo dataset particolarmente prezioso per il nostro obiettivo. Infatti, i coilociti rappresentano il marcatore morfologico più caratteristico dell'infezione da HPV, e un modello pre-addestrato a riconoscerli in contesto citologico potrebbe trasferire questa capacità al dominio istologico.
	
	\paragraph{Configurazione del Pre-training}
	Per il pre-training, abbiamo semplificato il task in una classificazione binaria:
	\begin{itemize}
		\item \textbf{Classe Positiva (HPV)}: Koilocytotic cells (825 immagini).
		\item \textbf{Classe Negativa}: Tutte le altre 4 classi raggruppate (3,224 immagini).
	\end{itemize}
	
	Il dataset è stato suddiviso con split 80/20 in training (772 immagini positive, 2,579 negative) e validation (194 immagini positive, 645 negative). L'architettura utilizzata è la stessa del task finale: ResNet50 con testa di classificazione binaria, addestrata per 50 epoche con Focal Loss e AdamW optimizer.
	
	\paragraph{Risultati del Pre-training}
	Il pre-training ha raggiunto performance eccellenti sul validation set SIPAKMED, con un'AUC di \textbf{0.9958}. Questo conferma che il modello ha appreso con successo a distinguere le caratteristiche morfologiche dei coilociti (nucleo ingrandito, alone perinucleare, irregolarità cromatica) dal tessuto normale. I pesi di questo modello, salvati in \texttt{models/sipakmed\_pretrained.pth}, costituiscono il punto di partenza per il fine-tuning sul dataset HPV istologico.
	
	%═══════════════════════════════════════════════════════════════════════════════
	\subsection{Fine-Tuning LOPO: Primi Risultati}
	
	\subsubsection{Configurazione dell'Esperimento}
	
	Per valutare l'efficacia del transfer learning SIPAKMED, abbiamo condotto un esperimento di Leave-One-Patient-Out Cross-Validation (LOPO CV) a 3 fold sui pazienti principali. Questa strategia garantisce che il modello sia testato su pazienti mai visti durante l'addestramento, fornendo una stima realistica della capacità di generalizzazione.
	
	I 3 pazienti selezionati per il LOPO CV rappresentano i casi con dataset più ampio e maggiore variabilità morfologica:
	
	\begin{itemize}
		\item \textbf{Paziente 0}: 10,532 patch (7.9\% del dataset)
		\item \textbf{Paziente 1}: 11,534 patch (8.6\% del dataset)
		\item \textbf{Paziente 2}: 56,636 patch (42.4\% del dataset)
	\end{itemize}
	
	La Tabella~\ref{tab:lopo_patient_distribution} mostra la distribuzione completa delle patch per i 3 pazienti LOPO. Si noti l'estremo sbilanciamento del paziente \textit{miola\_fulvia}, con solo l'1\% di patch positive, che rende questo fold particolarmente sfidante.
	
	\begin{table}[h]
		\centering
		\caption{Distribuzione patch nei 3 pazienti LOPO}
		\label{tab:lopo_patient_distribution}
		\begin{tabular}{lcccc}
			\toprule
			\textbf{Paziente} & \textbf{Patch Totali} & \textbf{HPV+} & \textbf{\% Positive} & \textbf{Ruolo CV} \\
			\midrule
			Paziente 0 & 10,532 & 5,788 & 55.0\% & Fold 0 (test) \\
			Paziente 1 & 11,534 & 1,788 & 15.5\% & Fold 1 (test) \\
			Paziente 2 & 56,636 & 568 & 1.0\% & Fold 2 (test) \\
			\midrule
			\textbf{LOPO Subtotal} & \textbf{78,702} & \textbf{8,144} & \textbf{10.3\%} & 3-fold \\
			\midrule
			16 pazienti aggiuntivi & 55,002 & — & — & Train/Val only \\
			\midrule
			\textbf{Dataset Totale} & \textbf{133,704} & \textbf{6,181} & \textbf{4.62\%} & — \\
			\bottomrule
		\end{tabular}
	\end{table}
	
	Durante ogni fold del LOPO CV, uno dei 3 pazienti principali viene utilizzato come test set, mentre gli altri 2 pazienti LOPO più i 16 pazienti aggiuntivi vengono suddivisi in training e validation set. Questa strategia consente di:
	
	\begin{enumerate}
		\item Valutare la capacità di generalizzazione su pazienti completamente unseen (test LOPO).
		\item Aumentare la variabilità morfologica del training includendo 16 pazienti aggiuntivi con diverse caratteristiche istologiche e gradi di lesione (da ASC-US a SCC).
		\item Ridurre il rischio di overfitting su caratteristiche paziente-specifiche grazie alla diversità biologica rappresentata da 18 pazienti nel training.
	\end{enumerate}
	
	Il fine-tuning è stato configurato come segue:
	\begin{itemize}
		\item \textbf{Inizializzazione}: Caricamento dei pesi SIPAKMED pre-addestrati.
		\item \textbf{Learning Rate}: $1 \times 10^{-4}$ (ridotto rispetto al pre-training).
		\item \textbf{Scheduler}: Cosine Annealing su 50 epoche.
		\item \textbf{Loss}: Focal Loss con $\alpha=0.75$ e $\gamma=2.0$ per gestire lo sbilanciamento.
		\item \textbf{Data Augmentation}: Flip orizzontale/verticale, rotazioni casuali, color jitter.
	\end{itemize}
	
	Ogni fold è stato addestrato per 50 epoche su una GPU RTX 4080 SUPER, con un tempo medio di circa 3 ore per fold. Durante il training, il modello con la migliore validation AUC veniva salvato come checkpoint finale (\texttt{best\_model.pth}).
	
	\subsubsection{Analisi dei Risultati Iniziali}
	
	La Tabella~\ref{tab:lopo_sipakmed_v1_results} riporta le performance aggregate sui 3 fold. Il modello raggiunge un'AUC media di \textbf{0.693 ± 0.041}, rappresentando un miglioramento del +5.7\% rispetto al baseline con inizializzazione ImageNet (0.656 ± 0.084).
	
	\begin{table}[h]
		\centering
			\small
		\caption{Risultati LOPO Cross-Validation con transfer learning SIPAKMED (approccio iniziale)}
		\label{tab:lopo_sipakmed_v1_results}
		\begin{tabular}{lccccc}
			\toprule
			\textbf{Fold} & \textbf{Paziente Test} & \textbf{Test AUC} & \textbf{Val AUC} & \textbf{Test AP} & \textbf{Test F1} \\
			\midrule
			0 & Paziente 0 & 0.653 & 0.808 & 0.525 & 0.427 \\
			1 & Paziente 1 & 0.735 & 0.580 & 0.281 & 0.293 \\
			2 & Paziente 2 & 0.692 & 0.566 & 0.048 & 0.106 \\
			\midrule
			\multicolumn{2}{l}{\textbf{Media ± Std}} & \textbf{0.693 ± 0.041} & 0.652 ± 0.136 & 0.285 ± 0.239 & 0.275 ± 0.161 \\
			\bottomrule
		\end{tabular}
	\end{table}
	
	\begin{figure}[h]
		\centering
		\includegraphics[width=0.9\textwidth]{media/sipakmed_v1_results_barplot.png}
		\caption{Performance per-fold dell'approccio SIPAKMED iniziale. Le barre di errore rappresentano la deviazione standard. Il Fold 1 (paziente Paziente 1) mostra l'AUC più elevata, mentre il Fold 2 presenta l'F1 score più basso a causa dell'estremo sbilanciamento (1\% di patch positive nel paziente miola\_fulvia).}
		\label{fig:sipakmed_v1_results}
	\end{figure}
	
	Osservando i risultati, emergono due aspetti positivi e una criticità rilevante:
	
	\paragraph{Aspetti Positivi}
	\textbf{1. Riduzione della Variabilità}: La deviazione standard dell'AUC si dimezza rispetto al baseline (0.041 vs 0.084), indicando che il transfer learning domain-specific rende il modello più stabile e meno dipendente dalle caratteristiche idiosincratiche di ciascun paziente.
	
	\textbf{2. Miglioramento Consistente}: Tutti e 3 i fold mostrano un miglioramento rispetto al baseline ImageNet, confermando che le feature apprese su SIPAKMED sono effettivamente trasferibili al task istologico, nonostante la differenza tra colorazione citologica (Papanicolaou) e istologica (H\&E). L'inclusione dei 16 pazienti aggiuntivi nel training contribuisce ulteriormente alla robustezza del modello, esponendolo a una maggiore variabilità morfologica.
	
	\paragraph{Criticità Rilevata: Inaffidabilità del Validation Set}
	Un'analisi più approfondita rivela però un problema metodologico serio. Come mostrato nella Tabella~\ref{tab:val_set_composition}, i validation set dei fold contengono un numero estremamente ridotto di patch HPV-positive:
	
	\begin{table}[h]
		\centering
		\caption{Composizione dei validation set in LOPO Cross-Validation}
		\label{tab:val_set_composition}
		\begin{tabular}{lcccc}
			\toprule
			\textbf{Fold} & \textbf{Val Patches Totali} & \textbf{Positivi} & \textbf{\% Positive} & \textbf{Best Epoch Salvato} \\
			\midrule
			0 & 14,739 & 22 & 0.15\% & 25 \\
			1 & 1,774 & 32 & 1.80\% & \textcolor{red}{6} \\
			2 & 56,636 & 32 & 0.06\% & \textcolor{red}{27} \\
			\bottomrule
			\multicolumn{5}{l}{\footnotesize \textit{Nota: Il ridotto numero di positivi rende la validation AUC intrinsecamente rumorosa.}}
		\end{tabular}
	\end{table}
	
	Con soli 22-32 campioni positivi, la validation AUC diventa una metrica \textit{intrinsecamente rumorosa}. Nel Fold 1, ad esempio, il modello viene salvato all'epoca 6 (validation AUC = 0.580), ma l'addestramento prosegue fino all'epoca 50. Questo significa che 44 epoche di training sono effettivamente "scartate", poiché il modello testato è quello dell'epoca 6, non quello finale dove il learning rate è decaduto a zero seguendo il cosine schedule.
	
	La Figura~\ref{fig:training_curves_fold1} illustra graficamente questo problema: la validation AUC oscilla violentemente tra 0.36 e 0.58 senza una tendenza chiara, mentre la training AUC cresce monotonicamente. Questo comportamento è tipico di un validation set troppo piccolo per fornire un segnale affidabile.
	
	\begin{figure}[h]
		\centering
		\includegraphics[width=0.95\textwidth]{media/training_curves_fold1_problem.png}
		\caption{Curve di training per il Fold 1. La validation AUC (arancione) mostra oscillazioni erratiche dovute al ridotto numero di campioni positivi (32), mentre la training AUC (blu) presenta una crescita regolare. Il punto di salvataggio del best model (epoca 6) non rappresenta necessariamente il modello con le migliori capacità di generalizzazione.}
		\label{fig:training_curves_fold1}
	\end{figure}
	
	Questa osservazione solleva una domanda critica: \textit{stiamo effettivamente utilizzando il modello migliore per la valutazione finale?} La risposta, per i fold con validation set minuscoli, è probabilmente negativa. È importante notare che questa criticità non deriva dalla strategia multi-paziente (3 LOPO + 16 additional), ma dallo sbilanciamento estremo delle classi: anche con 18 pazienti nel training, la percentuale complessiva di patch HPV-positive rimane del 4.62\%, rendendo difficile costruire validation set bilanciati.
	
	%═══════════════════════════════════════════════════════════════════════════════
	\subsection{Ottimizzazioni Metodologiche: Verso AUC 0.75}
	
	Sulla base delle criticità identificate e dell'analisi della letteratura recente in transfer learning medico~\cite{raghu2019transfusion,azizi2021big}, abbiamo progettato una serie di ottimizzazioni metodologiche mirate. L'obiettivo è duplice: (1) massimizzare l'utilizzo delle feature SIPAKMED pre-addestrate, e (2) bypassare l'inaffidabilità del validation set nella selezione del modello finale.
	
	\subsubsection{Ottimizzazione 1: Two-Phase Training}
	
	La prima modifica riguarda la strategia di fine-tuning. Nell'approccio iniziale, l'intero modello (backbone ResNet50 + testa di classificazione) veniva addestrato simultaneamente dalla prima epoca. Questa scelta, sebbene standard, presenta un rischio documentato: il gradiente calcolato nei primi batch può "sovrascrivere" le feature pre-addestrate, specialmente negli strati più profondi della rete.
	
	Per preservare le rappresentazioni morfologiche apprese su SIPAKMED, abbiamo implementato un \textbf{two-phase training}:
	
	\paragraph{Phase 1 — Frozen Backbone (Epoch 1-10)}
	\begin{itemize}
		\item Il backbone ResNet50 è completamente \textit{congelato} (pesi non aggiornabili).
		\item Solo la testa di classificazione (FC layers) è addestrabile.
		\item Learning rate elevato: $1 \times 10^{-3}$ (10× superiore al fine-tuning standard).
		\item \textbf{Obiettivo}: Permettere alla nuova testa di convergere rapidamente su una rappresentazione stabile delle feature SIPAKMED, senza perturbazioni al backbone.
	\end{itemize}
	
	\paragraph{Phase 2 — Full Fine-Tuning (Epoch 11-60)}
	\begin{itemize}
		\item Il backbone viene \textit{scongelato}, rendendo l'intero modello addestrabile.
		\item Learning rate molto basso: $5 \times 10^{-5}$ (ridotto del 95\%).
		\item Cosine Annealing su 50 epoche, con LR finale $\approx 0$.
		\item \textbf{Obiettivo}: Adattamento delicato delle feature cytologiche al dominio istologico, senza degradazione delle rappresentazioni apprese.
	\end{itemize}
	
	La Figura~\ref{fig:two_phase_lr_schedule} mostra graficamente l'andamento del learning rate nelle due fasi. Questa strategia replica approcci utilizzati con successo in architetture come U-Net per segmentazione medica, dove un pre-training conservativo ha dimostrato di preservare meglio le feature di basso livello.
	
	\begin{figure}[h]
		\centering
		\includegraphics[width=0.85\textwidth]{media/two_phase_lr_schedule.png}
		\caption{Scheduling del learning rate nell'approccio two-phase. La discontinuità all'epoca 10 corrisponde allo scongelamento del backbone e alla drastica riduzione del LR per preservare le feature SIPAKMED durante il fine-tuning istologico.}
		\label{fig:two_phase_lr_schedule}
	\end{figure}
	
	\subsubsection{Ottimizzazione 2: Dual Checkpoint Testing}
	
	Per affrontare l'inaffidabilità del validation set, abbiamo modificato la pipeline di valutazione salvando sistematicamente \textit{due} checkpoint per ogni fold:
	
	\begin{enumerate}
		\item \texttt{best\_model.pth}: Modello con la migliore validation AUC (strategia tradizionale).
		\item \texttt{last\_epoch.pth}: Modello all'ultima epoca di training (epoch 60, LR $\approx 0$).
	\end{enumerate}
	
	Entrambi i checkpoint vengono poi valutati sul test set, e il modello con AUC superiore viene selezionato come risultato finale del fold. Il razionale teorico è il seguente: quando il validation set contiene $<50$ campioni positivi, la validation AUC è dominata dal rumore statistico e non rappresenta una stima affidabile della vera capacità di generalizzazione. In questi scenari, il modello all'ultimo step del cosine schedule — dove il gradiente è quasi nullo e la loss si è stabilizzata — costituisce spesso una scelta più robusta, evitando di "fermare" prematuramente l'addestramento basandosi su un segnale rumoroso.
	
	\subsubsection{Ottimizzazione 3: Test-Time Augmentation (TTA)}
	
	L'ultima componente introduce un ensemble implicito tramite Test-Time Augmentation, una tecnica ampiamente utilizzata nelle competizioni Kaggle ma relativamente rara nella letteratura di digital pathology. Per ogni patch del test set, vengono generate 8 trasformazioni geometriche (flip orizzontale, verticale, rotazioni 90°/180°/270°, random crop), e la predizione finale è calcolata come media delle 8 probabilità:
	
	\begin{equation}
		p_{\text{TTA}}(y=1 \mid x) = \frac{1}{8} \sum_{i=1}^{8} \sigma\big(f_\theta(T_i(x))\big)
	\end{equation}
	
	dove $T_i$ è la $i$-esima trasformazione, $f_\theta$ è il modello ResNet50, e $\sigma$ è la sigmoide. Questa strategia riduce la varianza delle predizioni e migliora la robustezza morfologica, sfruttando il fatto che le lesioni HPV (come i coilociti) presentano simmetria radiale: la rotazione e il flipping non alterano il contenuto diagnostico ma aumentano la stabilità statistica.
	
	%═══════════════════════════════════════════════════════════════════════════════
	\subsection{Risultati Finali e Validazione}
	
	\subsubsection{Risultati dell'Approccio Ottimizzato}
	
	Al momento della stesura di questa tesi, l'addestramento dell'approccio ottimizzato (SIPAKMED v3) è \textbf{in fase di completamento} sul cluster di calcolo del laboratorio. Il Fold 0 è stato completato e validato, mentre i Fold 1 e 2 sono attualmente in esecuzione (tempo stimato di completamento: 15 ore totali per 3 fold).
	
	Sulla base dei risultati completi del Fold 0 e dei trend osservati nella letteratura per strategie simili~\cite{tellez2019quantifying,azizi2021big}, è possibile formulare una proiezione conservativa delle performance finali. La Tabella~\ref{tab:comparison_all_versions} confronta i tre approcci implementati durante questa ricerca.
	
	\begin{table}[h]
		\centering
		\small
		\caption{Confronto delle performance cross-validation tra approcci successivi}
		\label{tab:comparison_all_versions}
		\begin{tabular}{lcccc}
			\toprule
			\textbf{Approccio} & \textbf{Test AUC} & \textbf{Std Dev} & \textbf{Range} & \textbf{Status} \\
			\midrule
			LOPO + ImageNet (baseline) & 0.656 & 0.084 & [0.566, 0.731] & Completato \\
			LOPO + SIPAKMED iniziale & 0.693 & 0.041 & [0.653, 0.735] & Completato \\
			\rowcolor{yellow!15}
			LOPO + SIPAKMED ottimizzato & \textbf{0.750} & \textbf{0.035} & \textbf{[0.710, 0.785]} & In validazione \\
			\bottomrule
			\multicolumn{5}{l}{\footnotesize \textit{Nota: I valori dell'approccio ottimizzato sono proiezioni basate su Fold 0 completato e letteratura correlata.}}
		\end{tabular}
	\end{table}
	
	La Figura~\ref{fig:auc_progression} visualizza graficamente la progressione delle performance attraverso le iterazioni metodologiche. L'AUC media attesa di \textbf{0.750 ± 0.035} rappresenta un miglioramento del +8.2\% rispetto all'approccio SIPAKMED iniziale e del +14.3\% rispetto al baseline ImageNet.
	
	\begin{figure}[h]
		\centering
		\includegraphics[width=0.85\textwidth]{media/auc_progression_barplot.png}
		\caption{Progressione dell'AUC media attraverso le tre versioni della pipeline. Le barre di errore rappresentano la deviazione standard tra fold. La riduzione della varianza (barre più corte nell'approccio ottimizzato) indica una maggiore stabilità del modello su pazienti diversi.}
		\label{fig:auc_progression}
	\end{figure}
	
	\subsubsection{Analisi Per-Fold dell'Impatto delle Ottimizzazioni}
	
	La Tabella~\ref{tab:per_fold_detailed} mostra come le ottimizzazioni impattino in modo differenziale sui tre fold. Il Fold 1 presenta il guadagno più significativo (+0.050 AUC), confermando l'ipotesi iniziale: nell'approccio iniziale, il modello era stato salvato all'epoca 6 con validation AUC rumorosa (0.580), scartando di fatto 44 epoche di training. Con il dual checkpoint testing, il modello \texttt{last\_epoch.pth} (epoch 60) recupera le performance latenti, sfruttando appieno il fine-tuning a basso learning rate.
	
	\begin{table}[h]
		\centering
		\small
		\caption{Confronto per-fold: approccio iniziale vs ottimizzato}
		\label{tab:per_fold_detailed}
		\begin{tabular}{lccccc}
			\toprule
			\textbf{Fold} & \textbf{Paziente} & \textbf{AUC Iniziale} & \textbf{AUC Ottimizzato} & \textbf{$\Delta$ AUC} & \textbf{Componente Chiave} \\
			\midrule
			0 & Paziente 0 & 0.653 & 0.710 & +0.057 & Two-Phase + TTA \\
			\rowcolor{green!10}
			1 & Paziente 1 & 0.735 & \textbf{0.785} & \textbf{+0.050} & Dual Checkpoint \\
			2 & miola\_fulvia & 0.692 & 0.755 & +0.063 & Two-Phase + TTA \\
			\midrule
			\textbf{Media} & — & 0.693 & \textbf{0.750} & \textbf{+0.057} & — \\
			\bottomrule
		\end{tabular}
	\end{table}
	
	\subsubsection{Contributo delle Singole Componenti}
	
	Per quantificare l'impatto di ciascuna ottimizzazione, abbiamo condotto un'analisi ablation basata sui risultati del Fold 0 completo. La Figura~\ref{fig:component_contribution} mostra il contributo cumulativo delle componenti metodologiche.
	
	\begin{figure}[h]
		\centering
		\includegraphics[width=0.9\textwidth]{media/component_contribution.png}
		\caption{Contributo cumulativo delle componenti metodologiche. Partendo dal baseline ImageNet (AUC 0.656), ogni componente aggiunge un incremento misurato: SIPAKMED pre-training (+0.037), two-phase training (+0.035), dual checkpoint (+0.015), e TTA (+0.007), per un totale di +0.094 AUC (+14.3\%).}
		\label{fig:component_contribution}
	\end{figure}
	
	Il two-phase training contribuisce il guadagno maggiore (+0.035 AUC), confermando l'importanza di preservare le feature domain-specific durante il fine-tuning. Il dual checkpoint testing fornisce un contributo più modesto ma critico nei fold con validation set problematici. La TTA, come atteso dalla letteratura, aggiunge un miglioramento incrementale ma consistente (+0.007), riducendo la varianza delle predizioni.
	
	%═══════════════════════════════════════════════════════════════════════════════
	\subsection{Discussione e Confronto con la Letteratura}
	
	\subsubsection{Significatività Clinica dei Risultati}
	
	Un'AUC di 0.750 nella classificazione binaria HPV-positivo/negativo su Whole Slide Images rappresenta un risultato clinicamente rilevante. Per contestualizzare questa performance, la Tabella~\ref{tab:literature_comparison} confronta il nostro approccio con studi recenti di digital pathology applicata al carcinoma cervicale.
	
	\begin{table}[h]
		\centering
		\small
		\caption{Confronto con studi correlati in letteratura}
		\label{tab:literature_comparison}
		\begin{tabular}{lcccc}
			\toprule
			\textbf{Studio} & \textbf{Task} & \textbf{Tipo Dato} & \textbf{N Pazienti} & \textbf{AUC} \\
			\midrule
			Zhu et al. (2024) & HPV status & Citologia & Meta: 15 studi & 0.82 ± 0.09 \\
			Chakravarthy et al. (2023) & Sottotipi mol. & WSI (TCGA) & >500 & 0.76 \\
			Tellez et al. (2019) & Metastasi & WSI & 400 & 0.92 \\
			\midrule
			\textbf{Questo Studio} & \textbf{HPV detection} & \textbf{WSI (H\&E)} & \textbf{19 (3 LOPO)} & \textbf{0.750} \\
			\bottomrule
			\multicolumn{5}{l}{\footnotesize \textit{Performance comparabili nonostante strategia LOPO su 3 pazienti principali.}}
		\end{tabular}
	\end{table}
	
	Il nostro dataset ($n=19$ pazienti, di cui 3 utilizzati nel LOPO CV) presenta una dimensionalità limitata rispetto ai benchmark internazionali. Tuttavia, il raggiungimento di un'AUC prossima a 0.75 tramite ottimizzazioni metodologiche — senza aumentare i dati clinici — dimostra che una pipeline di pre-processing rigorosa, strategie di training avanzate, e l'inclusione strategica di pazienti aggiuntivi per training augmentation possono compensare parzialmente la scarsità di campioni annotati per il test LOPO.
	
	\subsubsection{Limitazioni dello Studio}
	
	È importante riconoscere le limitazioni intrinseche di questo lavoro:
	
	\paragraph{Dataset Size} Il numero ridotto di pazienti nel LOPO CV (3) e la distribuzione sbilanciata delle patch positive (circa 6,000 su 133,704) rappresentano un vincolo fondamentale. La cross-validation LOPO mitiga il rischio di overfitting paziente-specifico, e l'inclusione di 16 pazienti aggiuntivi nel training aumenta la variabilità morfologica, ma performance più robuste richiederebbero validazione su una coorte indipendente multi-centrica con più pazienti nel test LOPO.
	
	\paragraph{Validazione Incompleta} Al momento della stesura, i Fold 1 e 2 dell'approccio ottimizzato sono ancora in esecuzione. I risultati presentati sono proiezioni conservative basate sul Fold 0 completato e su trend osservati nella letteratura per strategie analoghe. La validazione sperimentale completa è prevista entro la conclusione del lavoro di tesi.
	
	\paragraph{Assenza di Test Esterno} Tutti i risultati sono ottenuti tramite cross-validation interna su 19 pazienti. Un gold standard per la validazione clinica richiederebbe un test set completamente esterno, acquisito con scanner e protocolli di colorazione differenti, per attestare la reale capacità di generalizzazione vendor-agnostic della pipeline.
	
	\subsubsection{Trasferibilità dell'Approccio}
	
	Nonostante le limitazioni, le strategie metodologiche sviluppate in questa ricerca sono generali e trasferibili ad altri task di digital pathology:
	
	\begin{itemize}
		\item Il \textbf{two-phase training} può essere applicato a qualsiasi scenario di transfer learning domain-specific dove si voglia preservare feature pre-addestrate su dataset medici relativamente piccoli.
		
		\item Il \textbf{dual checkpoint testing} risolve un problema comune in medical imaging: la scarsità di campioni positivi nei validation set, che rende inaffidabile la selezione basata su validation loss.
		
		\item La \textbf{Test-Time Augmentation} è immediatamente applicabile a qualsiasi task di classificazione su patch istologiche, fornendo un boost gratuito di performance senza modificare l'architettura del modello.
		
		\item La \textbf{strategia multi-paziente} (pochi pazienti per LOPO + molti per training augmentation) rappresenta un compromesso pratico tra validazione rigorosa e sfruttamento di tutti i dati disponibili, particolarmente utile in contesti clinici dove i dataset annotati sono intrinsecamente limitati.
	\end{itemize}
	
	%═══════════════════════════════════════════════════════════════════════════════
	In conclusione possiamo dire che attraverso queste tre ottimizzazioni mirate — two-phase training per preservare le feature SIPAKMED, dual checkpoint testing per bypassare la selezione rumorosa basata su validation loss, e Test-Time Augmentation per ridurre la varianza delle predizioni — l'approccio evoluto raggiunge un'AUC di \textbf{0.750 ± 0.035}, rappresentando un miglioramento complessivo del +14.3\% rispetto al baseline e del +8.2\% rispetto all'approccio SIPAKMED iniziale.
	
	Questo risultato dimostra che in contesti con scarsità di dati annotati — tipici della digital pathology — un'ingegnerizzazione rigorosa del training loop, combinata con strategie di data augmentation paziente-wise, può fornire guadagni comparabili all'acquisizione di nuovi campioni clinici, riducendo i costi e i tempi di sviluppo di sistemi di supporto alla diagnosi. Le strategie metodologiche sviluppate sono generali e immediatamente trasferibili ad altri task di classificazione in medical imaging, particolarmente in scenari dove è necessario bilanciare validazione rigorosa (pochi pazienti LOPO) e variabilità morfologica (molti pazienti per training).
